\chapter{Grundlagen: Die fraktale Dimension und die zwei Ebenen der Theorie}

Die Weber-de-Broglie-Bohm-Theorie mit fraktalem Raum (WDBT+) postuliert, dass der fundamentale Raum keine glatte Mannigfaltigkeit ist, sondern eine fraktale Struktur mit der Dimension

\begin{equation}
D = \frac{\ln(20\phi)}{\ln(2+\phi)} \approx 2,704, \qquad \phi = \frac{1+\sqrt{5}}{2}.
\end{equation}

Diese Dimension folgt aus:

\begin{itemize}
    \item Der Dodekaeder-Geometrie (20 Ecken),
    \item Dem Goldenen Schnitt \( \phi \),
    \item Einer Fibonacci-Rekursion \( N_{n+1} = 2N_n + N_{n-1} \).
\end{itemize}

Die Theorie ist zweistufig aufgebaut:

\begin{enumerate}
    \item \textbf{Informations-Weber-Theorie (IWT):} Fundamentale, diskrete Ur-Theorie. Information ist die einzige primäre Größe. Raum, Zeit, Materie und Energie emergieren aus der Struktur und Dynamik eines diskreten Informationsnetzwerks.
    \item \textbf{WDBT+:} Die kontinuumsnahe, emergente Physik, die aus der IWT im Grenzfall großer Skalen hervorgeht. Sie umfasst eine deterministische Quantentheorie, eine geschwindigkeitsabhängige Gravitation und eine stationäre Kosmologie.
\end{enumerate}

Die zentrale Aussage der WDBT+ ist: \emph{Der Raum ist nicht die Bühne der Physik – er ist die Projektion einer fraktalen Energieverteilung.}